%%%%%%%%%%%%%%%%%%%%%%%%%%%%%%%%%%%%%%%%%%%%%%%%%%%%%%%%%
%%%%%%%%%%%%%%%%%         FUNDAMENTOS              %%%%%%%%%%%%%%%%%%%%%%
%%%%%%%%%%%%%%%%%%%%%%%%%%%%%%%%%%%%%%%%%%%%%%%%%%%%%%%%%


\chapter{Objetivos}
\label{chap:objectives}

\chapterintro{
El desarrollo de esta Tesis Doctoral está motivado por la necesidad de avanzar en el campo de la bioinformática aplicada a la salud humana. En el contexto actual, la identificación de biomarcadores microbianos es fundamental para mejorar el diagnóstico y pronóstico de diversas enfermedades. Este capítulo detalla los objetivos que guían este trabajo de investigación, fundamentando la importancia de cada uno en el ámbito científico y médico.}


El objetivo principal de esta Tesis Doctoral es el \textbf{desarrollo y aplicación de metodologías bioinformáticas para la identificación de biomarcadores microbianos asociados a enfermedades humanas}. Esto se aborda con futuros fines diagnósticos, utilizando el análisis de datos de microbioma obtenidos mediante secuenciación masiva y técnicas de Machine Learning.


Por consiguiente, a partir del objetivo principal se plantean los siguientes objetivos específicos:
\begin{itemize}

    \item \textbf{Evaluar la capacidad predictiva y de estratificación del microbioma humano}. Este objetivo se centra en el desarrollo, implementación y optimización de flujos de trabajo bioinformáticos que permitan el análisis de datos de secuenciación masiva del microbioma. Se busca establecer la viabilidad computacional de utilizar las abundancias taxonómicas como variables de entrada para entrenar modelos de ML. La finalidad es construir modelos robustos capaces de clasificar y estratificar a los pacientes para fines diagnósticos (distinción entre individuos sanos y enfermos) o pronósticos (predicción de la respuesta o no respuesta a un determinado tratamiento), facilitando la identificación de subgrupos con perfiles microbianos distintivos y determinando el potencial del microbioma en la medicina personalizada.
    
    \item \textbf{Explorar el papel del microbioma humano como fuente de biomarcadores}. Una vez establecida la capacidad predictiva del microbioma, este objetivo se centra en aprovechar la información generada para el descubrimiento de biomarcadores. Se utilizarán modelos de ML interpretables para extraer los taxones que mayor contribución tengan en la clasificación o predicción de estados clínicos. La finalidad es caracterizar estos biomarcadores microbianos clave y su asociación con diversas patologías, utilizando la información generada como base para nuevas investigaciones.

    \item \textbf{Validar los biomarcadores microbianos identificados mediante un análisis comparativo en cohortes clínicas independientes}. Este objetivo se centra en evaluar la consistencia y robustez de los biomarcadores identificados inicialmente en una cohorte principal mediante comparaciones in silico con datos de cohortes externas. Esta validación externa garantiza que los biomarcadores mantengan su relevancia y potencial clínico en diferentes contextos, fortaleciendo su aplicabilidad en estudios futuros.
    
\end{itemize}